\section{Research scope and conventions}
\label{sec:scope}

The current ProveIt manuscript already contains extensive exact
polylogarithm, Gamma, Stieltjes, and correlation identities. Its newest
incoming reports go further: they construct arbitrary-depth ordered
Hurwitz germs, complete coincident Stieltjes products, and compute the
first two derivatives of Tornheim zeta on every admissible ray through
the origin \cite{Repo,IncomingOrdered,IncomingCoincident}. The questions
addressed here are specific continuations of that work.

\subsection{What is proved here}

\begin{enumerate}
\item The cubic digamma finite part is reduced to the strict harmonic
coefficient $\eta(1)$ already used in the ordered Hurwitz report:
\[
 \FP\int_0^1\psi(x)^3\dd x
       =3\zeta(2)-6\eta(1)-6\gamma\gamma_1.
\]
This identifies two previously separate coordinates in the inspected
reports. It supplies a convergent harmonic-series answer to the proposed
reduction of the diagonal Tornheim third derivative. It does not give a
finite evaluation of $\eta(1)$ using only ordinary zeta constants.
\item The entire cubic directional layer of $\T(A,B,C)$ is determined
by two fixed constants. An exact double-zeta product identity reduces
the number of unknown quadratic coefficients in its holomorphic
remainder; the resulting formula answers the asymmetric-ray question
in Section~9, Question~2 of the coincident report.
\item The mixed integer-order problem in Section~9, Question~11 of
that report has a finite constructive solution at arbitrary depth.
The free outer spectral order is retained throughout. The pole
description includes logarithmic singularities, which are absent
from the purely rational case.
\end{enumerate}

The Gaussian $S_6$ and revised $S_8$ candidates are different arithmetic
questions. They are not proved by the results below. The exact-identity
and source-status ledgers retain this distinction.

\subsection{Classical input and the extent of novelty}

The Mellin representation, Jonqui\`ere expansion, shuffle and stuffle
products, Mittag--Leffler methods, and the elementary relationship
between Hurwitz zeta and polygamma functions are classical
\cite{DLMFpolylog,DLMFhurwitz,DLMFgamma,Connon,BorweinDilcher}.
General Laurent algorithms for multiple zeta functions also predate
this report \cite{MOW}. Harmonic Stieltjes coefficients have an
established normalization and literature \cite{Kargin}.

The contribution is the specific cubic bridge, the complete formula
for arbitrary cubic directions and transverse curves, and the explicit
finite reduction and pole prescription for the mixed integer-order
sector. These are proved directly. They are new continuations relative
to the inspected repository corpus; a claim of global priority for
every equivalent formulation would require a broader literature audit.
In particular, a two-coordinate spanning formula is not a proof that
the corresponding numbers are arithmetically independent.

\subsection{Notation}

The generalized Stieltjes convention is
\begin{equation}\label{eq:stieltjes}
 \zeta(1+u,a)=\frac1u+
       \sum_{m=0}^{\infty}\frac{(-1)^m\gamma_m(a)}{m!}u^m,
 \qquad \gamma_m=\gamma_m(1),\quad \gamma=\gamma_0.
\end{equation}
Thus $\gamma_0(a)=-\psi(a)$ and $\psi=\Gamma'/\Gamma$.
Primes on $\zeta(s,a)$ denote derivatives in its first argument unless
explicitly stated otherwise. We set
\[
 L=\log(2\pi),\qquad
 H_n^{(r)}=\sum_{j=1}^n j^{-r},\qquad H_n=H_n^{(1)},\qquad H_0^{(r)}=0.
\]
Powers of positive real numbers use real logarithms. All logarithms
of complex variables in a right half-plane use the branch real on
the positive axis.

Our ordered double zeta convention places the larger summation index
first:
\begin{equation}\label{eq:Z2def}
 Z_2(s,t;a)=\sum_{n>m\ge0}(n+a)^{-s}(m+a)^{-t},\qquad a>0.
\end{equation}
The series converges, for example, if $\Re s>1$ and
$\Re(s+t)>2$. Its meromorphic continuation is understood when a
spectral argument is outside this domain. The strict harmonic coefficient
is fixed by
\begin{equation}\label{eq:eta-definition}
 Z_2(1+u,1;a)=\frac1{u^2}+\frac{\gamma_0(a)}u+
 \frac{\gamma_0(a)^2-\zeta(2,a)}2+\eta(a)u+O(u^2).
\end{equation}
This is the convention of the ordered incoming report. If
$\zeta_H(s,a)$ uses the inclusive inner sum, then
\begin{equation}\label{eq:strict-inclusive}
 \zeta_H(s,a)=Z_2(s,1;a)+\zeta(s+1,a),\qquad
 \eta(a)=-\gamma_H(1,a)-\zeta'(2,a).
\end{equation}
The correction $\zeta'(2,a)$ and the minus sign are essential.

For a function with an endpoint expansion in powers of $x$ and
$\log x$, $\FP\int_0^1$ means the constant coefficient of the cutoff
integral $\int_\varepsilon^1$ as $\varepsilon\downarrow0$, in the
coordinate $x$ itself. We never replace this prescription silently
by a spectral Laurent constant or by a nonlinear cutoff coordinate.
